\begin{problem}[Generic points and nonempty open subsets]
Let $X$ be irreducible with generic point $\eta$.
\begin{enumerate}
\item Prove that every nonempty open $U\subseteq X$ contains $\eta$.
\item Prove that $\eta$ is the generic point of $U$.
\item If $C$ is an irreducible component of a space with generic point $\eta_C$, show that
an open set $U$ meets $C$ if and only if $\eta_C\in U$.
\end{enumerate}
\end{problem}

\paragraph{Why this problem matters.}
It turns generic points into a practical test for whether an open set meets an irreducible
piece.

\paragraph{Useful ideas before starting.}
The complement of an open set is closed, and closures restrict to open subspaces.

\paragraph{Guided hints.}
If $\eta\notin U$, the closed complement contains all of $\overline{\{\eta\}}$.

\begin{solution}
If $U\neq\varnothing$ and $\eta\notin U$, then the closed set $X\setminus U$ contains
$\eta$.  Hence it contains
\[
\overline{\{\eta\}}=X,
\]
forcing $U=\varnothing$, contradiction.  Therefore $\eta\in U$.

In the subspace topology,
\[
\overline{\{\eta\}}^{\,U}
=
U\cap\overline{\{\eta\}}^{\,X}
=
U,
\]
so $\eta$ is generic in $U$.

For an irreducible component $C$, if $U\cap C\neq\varnothing$, then $U\cap C$ is a
nonempty open subset of $C$, so it contains $\eta_C$.  Conversely if $\eta_C\in U$, then
obviously $U\cap C\neq\varnothing$.
\end{solution}

\paragraph{What happened mathematically.}
The generic point is contained in exactly the open sets that genuinely meet its irreducible
component.

\paragraph{Extensions.}
Use this to characterize dense opens in a space with finitely many irreducible components.

\paragraph{Provenance.}
New problem based on the legacy ``generic point as general point'' discussion.
